\section{Verschränkung, Helizität und die Masselosigkeit des Photons}
\label{sec:verschränkung_masselosigkeit}

\subsection*{Überblick}
Ein oft übersehener, aber äußerst aufschlussreicher Hinweis auf die 
Masselosigkeit des Photons ergibt sich aus der Struktur seiner 
Verschränkungseigenschaften. Verschränkte Photonen verhalten sich 
experimentell exakt wie ein quantenmechanisches Zweizustands-System. 
Dies wiederum hängt direkt damit zusammen, dass ein masseloses 
Spin-1-Teilchen nur zwei physikalische Helizitätszustände besitzt. 
Schon eine beliebig kleine Photonenmasse würde diese Struktur verändern – 
ein Effekt, der experimentell klar nachweisbar wäre, aber nie beobachtet wurde.

\subsection*{Verschränkung überträgt keine physikalische Einwirkung}
Bei zwei verschränkten Photonen löst eine Messung an Photon~A 
\emph{keine} physikalische Veränderung an Photon~B aus. Es wird weder ein 
Signal noch Energie oder Information übertragen. Beide Photonen teilen 
sich einen gemeinsamen, nicht faktorisierbaren Quantenzustand. 
Eine Messung aktualisiert lediglich unser \emph{Wissen} über das 
Gesamtsystem, nicht den physikalischen Zustand des entfernten Photons. 
Verschränkung bedeutet daher keine Überlichtkommunikation und erfordert 
für sich genommen keine Masselosigkeit.

\subsection*{Helizitätsstruktur eines masselosen Spin-1-Teilchens}
Der entscheidende Punkt liegt woanders: Ein masseloses Spin-1-Teilchen 
wie das Photon besitzt nur zwei physikalische Freiheitsgrade,
\[
\lambda = +1, \qquad \lambda = -1,
\]
also die beiden transversalen Helizitätszustände.  
Diese Reduktion von drei auf zwei Zustände ist nur möglich, wenn
\[
m_\gamma = 0.
\]

Ein massives Spin-1-Teilchen besitzt hingegen stets drei Polarisationen,
\[
\lambda = +1, \qquad \lambda = 0, \qquad \lambda = -1,
\]
wobei der Zustand $\lambda = 0$ dem Longitudinalmodus entspricht.
Dieser Modus kann nur durch Eichsymmetrie entfernt werden – und eine 
Eichsymmetrie ist mit einem Proca-Masseterm unvereinbar.

\subsection*{Warum Verschränkungstests die Spinstruktur sichtbar machen}
Experimente mit verschränkten Photonen – darunter Bell-Tests, 
Polarisationskorrelationsmessungen, Hong--Ou--Mandel-Interferenz und 
quantentomografische Verfahren – reagieren äußerst empfindlich auf die 
Anzahl der zugänglichen Polarisationszustände.  
Alle beobachteten Verschränkungskorrelationen stimmen exakt mit der 
Struktur eines Zweizustands-Systems überein.

Hätte das Photon auch nur eine beliebig kleine Masse, so würde der 
Longitudinalmodus Beiträge liefern und damit messbar verändern:
\begin{itemize}
	\item die Winkelabhängigkeit von Bell-Korrelationen,
	\item die Sichtbarkeit von Polarisationsinterferenzen,
	\item die Struktur von Zwei-Photonen-Koinzidenzraten,
	\item das Verhalten unter Lorentztransformationen,
	\item die Rekonstruktion der Dichtematrix in der Quantentomografie.
\end{itemize}
Keine dieser Abweichungen wurde jemals beobachtet.

\subsection*{Fehlen eines Longitudinalmodus als Evidenz für Masselosigkeit}
Die empirische Tatsache, dass sämtliche Verschränkungsexperimente mit 
Photonen exakt mit zwei Polarisationszuständen konsistent sind, stellt 
eine starke indirekte Evidenz für die Masselosigkeit des Photons dar. 
Ein dritter Polarisationszustand wäre experimentell unübersehbar.
Sein eindeutiges Fehlen unterstützt:
\[
m_\gamma \approx 0
\quad\text{mit extrem strengen Grenzen aus Labor- und Astrophysikdaten}.
\]

\subsection*{Fazit}
Verschränkung vermittelt keine physikalische Einwirkung zwischen 
entfernten Photonen. Sie offenbart jedoch, welche Polarisationszustände 
ein Photon überhaupt besitzen kann. Die experimentell bestätigte 
Zweizustands-Struktur ist nur mit einem masselosen Spin-1-Feld 
vereinbar. Bereits eine minimale Photonenmasse würde einen dritten Modus 
aktivieren und die beobachteten Verschränkungssignaturen deutlich 
verändern.

Damit liefern die konsistenten Zweizustands-Eigenschaften verschränkter 
Photonen einen überzeugenden indirekten Hinweis darauf, dass das Photon 
mit extrem hoher Genauigkeit masselos ist.
